\section{M2 Mobile Capture Implementation}\label{sec:m2-mobile}

\noindent\fbox{\parbox{\dimexpr\linewidth-2\fboxsep-2\fboxrule\relax}{%
\textbf{Source code.}\\
Backend (Health service, in-scope API, JWT resource server):
\url{https://github.com/donga2606/vinabeef-m2}, commit
\texttt{99b87af} on \texttt{main}, 22 November 2026.\cite{m2BackendMobile}\\
Mobile client (Flutter):
\url{https://github.com/jaysonphan/vinabeef-mobile}, commit
\texttt{5d60fc7} on \texttt{main}, 22 November 2026.\cite{m2MobileRepo}\\
The Health backend verified in \Cref{sec:m2-health} is the earlier
\texttt{staging} commit \texttt{fee7a8d}.\cite{m2HealthStaging}}}

\medskip
This section covers the second M2 issue: a barn supervisor (BS) or farm
manager (FM) records a vaccination or health observation at the barn
with a phone and an RFID reader, while the barn may have no network.
Phan Ho Duy Thong implemented the mobile client and the server
additions it needs. The design is recorded in ADR-13 (\Cref{adr:13}),
the views in \Cref{sec:4.6}, and the requirement in FR-26.

\subsection{End-to-end flow}\label{sec:m2-mobile-flow}

\Cref{fig:15} shows the full flow. The numbered steps match the screens
of the app.

\begin{enumerate}
\item \textbf{Prepare session (online).} The user signs in with
Keycloak (Authorization Code with PKCE). The app calls
\texttt{GET /api/v1/cattle/in-scope} and caches the animals the user
may treat, with their \texttt{resourceVersion} and open vaccination
schedules. The cache is valid for 12 hours. Without a valid session
the app does not allow offline capture.
\item \textbf{Scan.} The reader produces a \texttt{TagObservation}
(source \texttt{MOBILE}, no direction). If several tags are read, the
user must pick one. An unknown or out-of-scope tag stays unresolved;
no animal is created.
\item \textbf{Confirm animal.} The app shows the matched animal (RFID,
farm, barn, open doses). The user must confirm it explicitly before a
form opens.
\item \textbf{Record at the barn.} The user fills in a vaccination
(vaccine, batch, dose, optional schedule) or a health observation. The
form is stored as a \texttt{DRAFT} in local SQLite and can still be
edited or discarded.
\item \textbf{Confirm the record.} Confirmation creates the
\texttt{intentId}, freezes the canonical request body (including
\texttt{expectedVersion}) and moves the intent to
\texttt{PENDING\_SYNC}. From here the content cannot change. The
screen says \emph{Saved on this device}, never \emph{Accepted}.
\item \textbf{Sync.} The sync coordinator runs on start, every 20 s,
after a confirmation and on demand. It posts the frozen body with the
same \texttt{intentId}. If an earlier response was lost, it first reads
\texttt{GET /health-intents/\{id\}} (the receipt) and only resubmits on
404. HTTP 401 keeps all intents pending and asks the user to sign in.
\item \textbf{Result.} The server outcome is stored and shown:
\texttt{ACCEPTED} (201 or 200 replay), \texttt{NEEDS\_REVIEW} (version
changed), or \texttt{REJECTED} (for example
\texttt{SCHEDULE\_ALREADY\_COMPLETED} or \texttt{ACCESS\_DENIED}). The
animal history screen reads the server history.
\end{enumerate}

The key distinction is between a \emph{local draft saved on the
device} and a \emph{record accepted by the server}. Only the Health
service accepts a clinical record, after it checks current
authorization, custody and version in one transaction (ADR-08).

\subsection{Module structure}\label{sec:m2-mobile-modules}

\Cref{fig:12} shows the source dependencies. The mobile client follows
the same ports-and-adapters split as \texttt{herd-health}.

{\small
\begin{longtable}{|p{0.27\linewidth}|p{0.64\linewidth}|}
\hline
\textbf{Package} & \textbf{Responsibility} \\ \hline
\endhead
\texttt{capture/domain} & \texttt{Cattle}, \texttt{LocalIntent},
\texttt{IntentStatus}, \texttt{SubmitResult}, \texttt{TagObservation},
\texttt{WorkSession}. No Flutter or I/O imports. \\ \hline
\texttt{capture/application} & Ports (\texttt{DraftStore},
\texttt{HealthApi}, \texttt{TagReadSource}, \texttt{WorkSessionStore},
\texttt{AuthSession}); \texttt{CaptureService} (validate, confirm,
freeze body); \texttt{SyncCoordinator} (receipt-first retry);
\texttt{WorkSessionService}. \\ \hline
\texttt{capture/data} & SQLite draft and session stores,
\texttt{HttpHealthApi} (dio), \texttt{FakeTagReadSource}, and an
in-app simulator that mirrors the Health contract for demos. \\ \hline
\texttt{capture/presentation} & Home, scan, confirm, form, intent
detail and history screens; status chips. \\ \hline
\texttt{auth} & \texttt{KeycloakAuthSession} (live) and
\texttt{DemoAuthSession}; refresh token in secure storage. \\ \hline
\texttt{composition\_root} & The only place that chooses demo or live
adapters. \\ \hline
\end{longtable}
}

\subsection{Server additions}\label{sec:m2-mobile-server}

The mobile flow needed two changes in \texttt{herd-health}, both in the
existing hexagonal structure:

\begin{itemize}
\item \texttt{GET /api/v1/cattle/in-scope}:
\texttt{CattleScopeController} $\rightarrow$
\texttt{ListCattleInScopeUseCase} $\rightarrow$
\texttt{CattleScopeService} $\rightarrow$ \texttt{CattleScopeQueryPort}
(\texttt{JdbcCattleScopeQuery}). The service pre-filters by the farms
in the caller's grants and checks \texttt{VIEW\_CATTLE} for each row,
so the list obeys the same policy as a single lookup. It returns
active animals with \texttt{resourceVersion} and open schedules. The
OpenAPI contract was updated.
\item JWT resource server: \texttt{HealthSecurityConfiguration}
validates Keycloak tokens (JWKS signature and issuer), maps the
principal from \texttt{preferred\_username}, permits only
\texttt{/actuator/health} without a token, and returns a
\texttt{problem+json} 401 for missing or invalid tokens. Current grants
still come from the grant adapter, which is a JSON fixture in this
prototype.
\end{itemize}

\subsection{Verification}\label{sec:m2-mobile-verification}

{\small
\begin{longtable}{|p{0.30\linewidth}|p{0.15\linewidth}|p{0.46\linewidth}|}
\hline
\textbf{Item} & \textbf{Status} & \textbf{Evidence} \\ \hline
\endhead
Backend build, all modules & Passed & \texttt{./dev mvn -B verify}, 22
Nov 2026: BUILD SUCCESS (access-policy 76, farm-edge 10, movement 3,
reader-emulator 14). Herd-health ran 20 tests; its 13 Testcontainers
tests are skipped in the toolbox (no Docker socket) and are covered by
\Cref{tab:m2-health-tests}. \\ \hline
\texttt{CattleScopeService} unit tests & Passed & 5 tests in the run
above. \\ \hline
Local stack (Compose, profile \texttt{local}) & Passed & All containers
healthy; \texttt{/actuator/health} returns UP. \\ \hline
Token required & Passed & \texttt{in-scope} without a token returns 401
\texttt{problem+json}. \\ \hline
Scope with real Keycloak token & Passed & Token for \texttt{bs.demo}
(BS, BARN-DL-01) returns exactly one animal, RFID-DL-1001, with the
open FMD-BOOSTER schedule; animals in BARN-DL-02 and FARM-GL are not
returned. \\ \hline
App on iOS simulator & Run & Demo mode screens exercised manually. \\ \hline
Mobile automated tests (24) & Not recorded & \texttt{flutter test};
the list is below. Results are not recorded in this report. \\ \hline
QAS09 fixture (20 drafts, 120 s) & Not run & Planned measure; see
\Cref{sec:qas09}. \\ \hline
\end{longtable}
}

The mobile tests cover the behaviour the review asked for:

\begin{itemize}
\item \textbf{Capture} (5): confirmation freezes the body with a stable
\texttt{intentId}; a submitted intent cannot be edited, re-confirmed
or deleted; invalid input and a future \texttt{occurredAt} are rejected
before an intent exists.
\item \textbf{Sync} (7): offline intents stay \texttt{PENDING\_SYNC}
and are later accepted once with the same id and body; a lost response
is resolved by the receipt; receipt 404 leads to a resubmit of the same
body; 401 keeps everything pending; final outcomes are not resent.
\item \textbf{Simulated Health semantics} (11): retry $\times$3 gives one
record; changed body gives \texttt{IDEMPOTENCY\_KEY\_REUSED}; completed
schedule gives \texttt{REJECTED}; a second offline intent for the same
animal gives \texttt{NEEDS\_REVIEW}; wrong barn and moved custody are
denied; unknown fields are rejected.
\item \textbf{Widget} (1): the status chip distinguishes pending from
accepted.
\end{itemize}

Demo scenarios mapped to the M2 cases: T05 (offline capture, then
sync), T06 (retry and version), T07 (wrong barn, revoked grant) and
T09 (unknown tag, explicit confirmation). In demo mode they run against
the in-app simulator; in live mode against the local Compose stack.

\subsection{How to run}\label{sec:m2-mobile-run}

\begin{lstlisting}[basicstyle=\footnotesize\ttfamily,breaklines=true,columns=fullflexible]
# backend (repository vinabeef-m2)
cp -n .env.example .env
./dev up -Profile local          # Postgres, Keycloak, herd-health ...
# mobile (repository vinabeef-mobile)
./scripts/bootstrap.sh ios       # or android
flutter test
flutter run                      # choose Demo or Live on the login screen
\end{lstlisting}

Live mode uses \url{http://localhost:8080} (Health) and
\url{http://localhost:8081/realms/vinabeef} (Keycloak). Demo users:
\texttt{bs.demo} and \texttt{fm.demo}.

\subsection{Limitations and open points}\label{sec:m2-mobile-limits}

\begin{itemize}
\item The RFID reader is simulated (\texttt{FakeTagReadSource}); the
Bluetooth reader adapter is not built.
\item Grants come from a JSON fixture, not a real IAM grant service.
\item Custody moved out of scope before sync returns 403 and the app
shows \texttt{REJECTED}. Treating it as \texttt{NEEDS\_REVIEW} is still
an open team decision.
\item Because \texttt{resourceVersion} is an optimistic lock, a second
offline record for the same animal before sync goes to
\texttt{NEEDS\_REVIEW}. This is the intended safe outcome, but it
creates review work for the FM.
\item The snapshot checker has not yet been run on a mobile export.
\end{itemize}
